% Preamble of the book: fonts, the listing style for Algol 68, and the
% macros of the text.  Everything here must also be understood by the HTML
% conversion (pandoc, see html/algol68.lua): plain \newcommand macros,
% listings, tables; no TikZ or other drawing.

\usepackage{fontspec}
\setmonofont{DejaVu Sans Mono}[Scale=MatchLowercase]
\usepackage[english]{babel}
\usepackage{microtype}
\setlength{\emergencystretch}{3em}
% Subsections (\subsection*), the headings within a section: in the size of
% the text, with at least one empty line above them.
\RedeclareSectionCommand[
  font=\normalsize\sffamily\bfseries,
  beforeskip=1.5\baselineskip plus 0.5\baselineskip minus 0.2\baselineskip,
  afterskip=0.5\baselineskip]{subsection}
\usepackage{xcolor}
\definecolor{LinkBlue}{rgb}{0.10,0.25,0.55}
\usepackage{booktabs}
\usepackage{longtable}
\usepackage{makeidx}
\makeindex
\usepackage[colorlinks,linkcolor=LinkBlue,urlcolor=LinkBlue,citecolor=LinkBlue,linktoc=all,bookmarksnumbered,pdfusetitle]{hyperref}

\usepackage{listings}
% A heading sets the flag @nobreak, which the first paragraph after it
% clears.  A listing directly after a heading empties \everypar, so the
% flag would stay set, and the next heading would omit its space above as
% if it followed the first one directly.  The end of a listing clears it.
\makeatletter
\lst@AddToHook{DeInit}{\global\@nobreakfalse}
\makeatother
% Algol 68 in SUPPER stropping: the reserved words of ga68 are printed in
% bold, as in the publication language of the Report.  Type names that a
% program declares (Node) and operators (ABS, a user-defined MAX) are
% printed as they are written.
\lstdefinelanguage{Algol68}{
  sensitive=true,
  morekeywords={access,andth,assert,at,begin,bits,bool,by,bytes,case,channel,char,
    compl,def,do,egg,elif,else,empty,end,esac,exit,false,fed,fi,file,flex,
    for,format,from,go,goto,heap,if,in,int,is,isnt,loc,long,mode,module,
    nest,nil,od,of,op,orel,ouse,out,par,postlude,prio,proc,pub,real,ref,
    sema,short,skip,string,struct,then,to,true,union,void,while,word},
  morecomment=[n]{\{}{\}},
  morestring=[d]{"},
}
\lstset{
  language=Algol68,
  basicstyle=\ttfamily\small,
  keywordstyle=\bfseries,
  commentstyle=\itshape\color{black!70},
  stringstyle=\color{black},
  columns=fullflexible,
  keepspaces=true,
  showstringspaces=false,
  upquote=true,
  breaklines=true,
  breakatwhitespace=true,
  xleftmargin=1.5em,
  aboveskip=0.6\baselineskip,
  belowskip=0.6\baselineskip,
}
\lstdefinestyle{output}{
  language={},
  basicstyle=\ttfamily\small\color{black!80},
  frame=leftline,
  framerule=1.5pt,
  rulecolor=\color{black!30},
  xleftmargin=1.5em,
  framexleftmargin=0.6em,
}
\lstdefinestyle{shell}{
  language={},
  basicstyle=\ttfamily\small,
  xleftmargin=1.5em,
}

% \example{ch02/first}: the program book/examples/ch02/first.a68 (its
% excerpt, see tools/run_examples.py) followed by its recorded output.
% \program{...}: the program only.  \outputof{...}: the output only.
% \programpart{ch07/json}{types}: the excerpt named types of the program.
\newcommand{\program}[1]{\lstinputlisting[language=Algol68]{build/listings/#1.a68}}
\newcommand{\programpart}[2]{\lstinputlisting[language=Algol68]{build/listings/#1.#2.a68}}
\newcommand{\cprogram}[1]{\lstinputlisting[language=C,basicstyle=\ttfamily\small]{examples/#1.c}}
\newcommand{\outputof}[1]{\lstinputlisting[style=output]{examples/#1.out}}
\newcommand{\example}[1]{\program{#1}\outputof{#1}}

% Inline Algol 68: \code|x := 1|.
\newcommand{\code}{\lstinline}

% \term{reference}: a term introduced here (italic, in the index).
% \rr{name}: the term the Revised Report uses for the same thing (in the
% index, so that a reader of the Report finds it).
\newcommand{\term}[1]{\emph{#1}\index{#1}}
\newcommand{\rr}[1]{\textsf{#1}\index{#1@\textsf{#1} (Report)}}

% A note on the implementations (ga68, a68j, Algol 68 Genie).
\newenvironment{implnote}{\begin{quote}\small\textbf{Implementation note.}\ }{\end{quote}}
\newenvironment{pitfall}{\begin{quote}\small\textbf{Pitfall.}\ }{\end{quote}}
% Room for three-digit page numbers in the table of contents
\DeclareTOCStyleEntries[pagenumberwidth=2.2em]{tocline}{chapter,section}
